% ============================================================================
%  Chapter 6 --- AI Models
% ============================================================================
\chapter{AI Models and Intelligence Layer}
\label{ch:ai}

\section{Introduction}

The intelligence layer of the AI Operations Agent comprises three machine learning models and one statistical correlation engine. Each addresses a distinct analytical need within the CEM--CVM bridge:

\begin{enumerate}
  \item \textbf{SLA Risk Prediction} --- supervised regression on aggregated OSS features.
  \item \textbf{Network (OSS) Anomaly Detection} --- unsupervised outlier detection on per-record KPIs.
  \item \textbf{Revenue (BSS) Anomaly Detection} --- unsupervised outlier detection on subscriber metrics.
  \item \textbf{OSS--BSS Correlation Engine} --- statistical analysis of metric pairs across domains.
\end{enumerate}

This chapter documents each model's design rationale, feature space, training methodology, and evaluation strategy.

\section{Model Overview}

% ── Figure: ML Model Pipeline ─────────────────────────────────────────────────
\begin{figure}[H]
\centering
\begin{tikzpicture}[node distance=1.5cm and 1.8cm]

  % Input data
  \node[service=cemteal, minimum width=3cm] (ossdata) {OSS Data\\{\scriptsize (real TT KPIs)}};
  \node[service=bssgreen, minimum width=3cm, below=1cm of ossdata] (bssdata) {BSS Data\\{\scriptsize (real TT subscribers)}};
  
  % Feature engineering
  \node[service=ossblue, right=2cm of ossdata, minimum width=3.5cm, yshift=-0.5cm] (feat) {
    \textbf{Feature}\\
    \textbf{Engineering}\\
    {\scriptsize Aggregate / Enrich}
  };
  
  % Models
  \node[service=aiviolet, right=2cm of feat, yshift=1.5cm, minimum width=3.5cm] (gbr) {
    \textbf{GBR}\\
    {\scriptsize SLA Risk Score}\\
    {\scriptsize $\in [0, 1]$}
  };
  \node[service=aiviolet, right=2cm of feat, yshift=0cm, minimum width=3.5cm] (ifoss) {
    \textbf{Isolation Forest}\\
    {\scriptsize OSS Anomalies}
  };
  \node[service=aiviolet, right=2cm of feat, yshift=-1.5cm, minimum width=3.5cm] (ifbss) {
    \textbf{Isolation Forest}\\
    {\scriptsize BSS Anomalies}
  };
  
  % Correlation
  \node[service=cvmorange, below=0.5cm of ifbss, minimum width=3.5cm] (corr) {
    \textbf{Correlation Engine}\\
    {\scriptsize Pearson + Spearman}
  };
  
  % Outputs
  \node[service=bssgreen, right=2cm of ifoss, minimum width=3cm, minimum height=3cm] (output) {
    \textbf{CVM-Ready}\\
    \textbf{Intelligence}\\[3pt]
    {\scriptsize Risk scores}\\
    {\scriptsize Anomaly alerts}\\
    {\scriptsize Correlation insights}
  };
  
  % Arrows
  \draw[myarrow=cemteal!70] (ossdata) -| (feat);
  \draw[myarrow=bssgreen!70] (bssdata) -| (feat);
  \draw[myarrow=ossblue!70] (feat) |- (gbr);
  \draw[myarrow=ossblue!70] (feat) -- (ifoss);
  \draw[myarrow=ossblue!70] (feat) |- (ifbss);
  \draw[myarrow=ossblue!70] (feat) |- (corr);
  \draw[myarrow=aiviolet!70] (gbr) -| (output);
  \draw[myarrow=aiviolet!70] (ifoss) -- (output);
  \draw[myarrow=aiviolet!70] (ifbss) -| (output);
  \draw[myarrow=cvmorange!70] (corr) -| (output);
  
\end{tikzpicture}
\caption{ML model pipeline: from raw data through feature engineering to CVM-ready intelligence.}
\label{fig:ml-pipeline}
\end{figure}

\begin{table}[H]
\centering
\caption{Model inventory: three ML models and one statistical engine.}
\label{tab:model-overview}
\begin{tabularx}{\textwidth}{@{}llllX@{}}
\toprule
\textbf{Model} & \textbf{Algorithm} & \textbf{Type} & \textbf{Features} & \textbf{Output} \\
\midrule
SLA Risk Scorer & GradientBoostingRegressor & Supervised regression & 9 aggregate OSS & Risk score $\in [0, 1]$ + importances \\
OSS Anomaly & IsolationForest & Unsupervised & 5 per-record OSS & Anomaly flag + severity $\in [0, 1]$ \\
BSS Anomaly & IsolationForest & Unsupervised & 7 per-record BSS & Anomaly flag + severity $\in [0, 1]$ \\
Correlation & Pearson + Spearman & Statistical & 5 OSS--BSS pairs & 10 correlation coefficients + p-values \\
\bottomrule
\end{tabularx}
\end{table}

\section{Model 1: SLA Risk Prediction (GradientBoostingRegressor)}

\subsection{Business Objective}

Predict the risk that a current KPI time window will lead to an SLA (Service Level Agreement) breach. The output is a continuous score in $[0.0, 1.0]$, where 0 indicates nominal performance and 1 indicates imminent SLA violation. This score maps directly to SmartCare's experience risk concept and feeds CVM retention decision-making.

\subsection{Why Gradient Boosting}

Table~\ref{tab:gbr-justification} justifies the model choice:

\begin{table}[H]
\centering
\caption{Justification for GradientBoostingRegressor.}
\label{tab:gbr-justification}
\begin{tabularx}{\textwidth}{@{}lX@{}}
\toprule
\textbf{Criterion} & \textbf{Rationale} \\
\midrule
Data type & Tabular, numeric --- GBR's strongest domain \\
Non-linearity & SLA risk has threshold-driven, non-linear dynamics (e.g., latency only matters after a threshold) \\
Feature interactions & High latency + packet loss is worse than either alone --- GBR captures this \\
Dataset size & Small-to-medium --- GBR excels; neural networks would overfit \\
Explainability & \texttt{feature\_importances\_} enables direct explanation during defence \\
Maturity & Well-understood, battle-tested, academically defensive \\
\bottomrule
\end{tabularx}
\end{table}

\subsection{Input Features (9 Aggregate OSS KPIs)}

The model takes as input an aggregated KPI feature vector per time window per region:

$$\mathbf{x} = \begin{bmatrix} \bar{T} & \sigma_T & \bar{L} & \sigma_L & L_{\max} & \bar{P} & P_{\max} & \bar{U} & \bar{S} \end{bmatrix}$$

Where $T$ = throughput (Mbps), $L$ = latency (ms), $P$ = packet loss (\%), $U$ = active users, $S$ = RSRP (dBm). Bars denote means, $\sigma$ denotes standard deviations.

\subsection{Target Variable}

The risk label $y \in [0, 1]$ is constructed from the features using a domain-knowledge-driven formula:

\begin{align}
y &= \text{clip}\!\left(\frac{\bar{L} - 20}{60},\; 0,\; 0.35\right) 
   + \text{clip}\!\left(\frac{L_{\max} - 30}{70},\; 0,\; 0.25\right) \nonumber\\
  &+ \text{clip}\!\left(\frac{\bar{P}}{4},\; 0,\; 0.25\right)
   + \text{clip}\!\left(\frac{P_{\max}}{6},\; 0,\; 0.15\right) \nonumber\\
  &+ \text{clip}\!\left(\frac{60 - \bar{T}}{100},\; 0,\; 0.20\right)
   + \epsilon, \quad \epsilon \sim \mathcal{N}(0, 0.03)
\label{eq:risk-label}
\end{align}

This formula encodes telecom domain knowledge: latency contributes risk only above a threshold (20~ms for mean, 30~ms for max), packet loss always contributes proportionally, low throughput (below 60~Mbps) adds risk, and a small noise term $\epsilon$ prevents the model from learning a trivially deterministic function.

With real TT data, this label construction is applied to the actual observed KPI values --- the model learns from \emph{real} degradation patterns, not synthetic injections.

\subsection{Training Pipeline}

\begin{lstlisting}[language=Python, caption={SLA Risk model training pipeline.}]
Pipeline([
    ("scaler", StandardScaler()),
    ("gbr", GradientBoostingRegressor(
        n_estimators=200,
        max_depth=4,
        learning_rate=0.05,
        subsample=0.8,
        random_state=42,
    )),
])
\end{lstlisting}

\subsection{Hyperparameters}

\begin{table}[H]
\centering
\caption{GBR hyperparameters and their roles.}
\label{tab:gbr-hyperparams}
\begin{tabularx}{\textwidth}{@{}llcX@{}}
\toprule
\textbf{Hyperparameter} & \textbf{Value} & \textbf{Tunable} & \textbf{Effect} \\
\midrule
\texttt{n\_estimators}  & 200  & Yes & Number of boosting stages \\
\texttt{max\_depth}     & 4    & Yes & Depth of each weak tree (complexity control) \\
\texttt{learning\_rate} & 0.05 & Yes & Contribution of each tree (lower = more stable) \\
\texttt{subsample}      & 0.8  & Yes & Stochastic sampling per stage (regularisation) \\
\texttt{random\_state}  & 42   & No  & Reproducibility seed \\
\bottomrule
\end{tabularx}
\end{table}

\subsection{Feature Importances and Explainability}

GBR provides \texttt{feature\_importances\_} --- the normalised total reduction of the loss function contributed by each feature. In practice, \texttt{mean\_latency\_ms} dominates (importance $\sim 0.69$), consistent with the risk label construction and real telecom behaviour: latency is the primary driver of subscriber experience degradation.

The feature importances are returned as part of every SLA risk inference response, enabling CVM systems (or human analysts) to understand \emph{why} the risk score is high or low.

\subsection{Evaluation Metrics}

For the SLA risk model (supervised regression):
\begin{itemize}
  \item \textbf{MAE} (Mean Absolute Error): Average magnitude of prediction errors, in risk-score units.
  \item \textbf{RMSE} (Root Mean Squared Error): Penalises larger errors more heavily.
  \item \textbf{$R^2$} (Coefficient of Determination): Proportion of variance explained by the model.
\end{itemize}

With real data, we use a train/validation/test split to obtain unbiased estimates.

\section{Model 2: OSS Anomaly Detection (IsolationForest)}

\subsection{Business Objective}

Detect abnormal network KPI records --- throughput collapses, latency spikes, packet loss surges, cell overloads, or weak signal conditions --- \emph{without} requiring labelled data. This corresponds to SmartCare's demarcation function: identifying \emph{which cells} are experiencing abnormal behaviour.

\subsection{Why Isolation Forest}

\begin{itemize}
  \item \textbf{No labels required:} In real telecom operations, comprehensive anomaly labels are rarely available. IsolationForest learns the ``normal'' distribution and flags deviations.
  \item \textbf{Speed and scalability:} Linear time complexity $O(n \cdot t \cdot \psi)$ where $t$ = trees, $\psi$ = subsample size.
  \item \textbf{Dual output:} Provides both a binary decision (normal/anomaly) and a continuous anomaly score for severity ranking.
  \item \textbf{Interpretable:} Anomaly score = normalised inverse of average path length. Shorter paths = more anomalous.
\end{itemize}

\subsection{Input Features (5 Per-Record OSS KPIs)}

$$\mathbf{x} = \begin{bmatrix} T & L & P & U & S \end{bmatrix}$$

Where $T$ = throughput, $L$ = latency, $P$ = packet loss, $U$ = active users, $S$ = RSRP. Each record represents one cell at one point in time.

\subsection{Anomaly Score Normalisation}

The raw IsolationForest \texttt{decision\_function} output is more negative for more anomalous records. We normalise to $[0, 1]$ where 1 = most anomalous:

$$\text{severity}_i = 1 - \frac{s_i - s_{\min}}{s_{\max} - s_{\min} + \epsilon}$$

where $s_i$ is the raw score and $\epsilon = 10^{-9}$ prevents division by zero.

\subsection{Contamination Parameter}

The \texttt{contamination} parameter defines the expected proportion of anomalies in the dataset. This is the \emph{most sensitive} hyperparameter in IsolationForest:

\begin{itemize}
  \item In real operator data, anomalies are typically rarer than the 5\% used in synthetic training. We anticipate contamination values in the range $[0.01, 0.03]$ for real TT data.
  \item The exact value will be calibrated during the retraining phase, using validation against known network events or operator-labelled incidents.
\end{itemize}

\subsection{Training Pipeline}

\begin{lstlisting}[language=Python, caption={OSS Anomaly model training pipeline.}]
Pipeline([
    ("scaler", StandardScaler()),
    ("ifo", IsolationForest(
        n_estimators=150,
        contamination=0.05,  # adjusted for real data
        random_state=42,
    )),
])
\end{lstlisting}

\subsection{Handling Class Imbalance in Real Data}

Real anomalies in production telecom networks are significantly rarer than the 5\% synthetic injection rate. Strategies for handling this:

\begin{enumerate}
  \item \textbf{Lower contamination:} Reduce to $0.01$--$0.03$ to match real prevalence.
  \item \textbf{Sliding threshold:} Instead of using the built-in decision threshold, apply a percentile-based threshold on anomaly scores.
  \item \textbf{Operator-guided calibration:} If TT provides a list of known incident dates/cells, we calibrate the threshold to maximise detection of those known events.
\end{enumerate}

\section{Model 3: BSS Revenue Anomaly Detection (IsolationForest)}

\subsection{Business Objective}

Detect anomalous subscriber-level business behaviour:
\begin{itemize}
  \item Abnormal recharge patterns (very high = potential SIM box fraud; very low = dormant SIM)
  \item Excessive data consumption outside normal plan bounds
  \item SMS spam patterns (very high SMS count)
  \item SIM-box-like voice patterns (very high voice minutes)
  \item Abrupt churn-risk spikes correlated with network events
\end{itemize}

This model is the BSS-side complement to the OSS anomaly detector. Together, they provide the AI Agent with a 360° view of both network and business anomalies.

\subsection{Input Features (7 Per-Record BSS Metrics)}

$$\mathbf{x} = \begin{bmatrix} R & D & V & S & C & A & G \end{bmatrix}$$

Where $R$ = revenue\_tnd, $D$ = data\_used\_gb, $V$ = voice\_min, $S$ = sms\_count, $C$ = churn\_risk, $A$ = appu\_tnd, $G$ = dou\_gb.

The inclusion of APPU and DOU (supervisor-directed) expands the feature space from 5 to 7, providing richer subscriber characterisation for anomaly detection.

\subsection{Training Pipeline}

\begin{lstlisting}[language=Python, caption={BSS Revenue Anomaly model training pipeline.}]
Pipeline([
    ("scaler", StandardScaler()),
    ("ifo", IsolationForest(
        n_estimators=150,
        contamination=0.05,  # adjusted for real data
        random_state=42,
    )),
])
\end{lstlisting}

\section{Model 4: OSS--BSS Correlation Engine}

\subsection{Objective}

Quantify the statistical relationship between network degradation and business impact. This is the \textbf{O+B convergence demonstration}: showing, with statistical evidence, that when the network degrades, revenue suffers.

\subsection{Metric Pairs}

The engine computes correlations on five OSS--BSS metric pairs:

\begin{table}[H]
\centering
\caption{OSS--BSS correlation pairs.}
\label{tab:corr-pairs}
\begin{tabular}{@{}lll@{}}
\toprule
\textbf{OSS Metric ($x$)} & \textbf{BSS Metric ($y$)} & \textbf{Expected Direction} \\
\midrule
\texttt{throughput\_mbps}    & \texttt{revenue\_tnd}   & Positive (higher throughput $\to$ higher revenue) \\
\texttt{latency\_ms}         & \texttt{revenue\_tnd}   & Negative (higher latency $\to$ lower revenue) \\
\texttt{packet\_loss\_pct}   & \texttt{churn\_risk}    & Positive (more loss $\to$ higher churn) \\
\texttt{latency\_ms}         & \texttt{data\_used\_gb} & Negative (higher latency $\to$ less data use) \\
\texttt{signal\_rsrp\_dbm}   & \texttt{revenue\_tnd}   & Positive (better signal $\to$ higher revenue) \\
\bottomrule
\end{tabular}
\end{table}

\subsection{Statistical Methods}

For each pair, we compute:

\begin{itemize}
  \item \textbf{Pearson $r$:} Measures linear correlation:
  $$r = \frac{\sum_{i=1}^{n}(x_i - \bar{x})(y_i - \bar{y})}{\sqrt{\sum(x_i - \bar{x})^2 \cdot \sum(y_i - \bar{y})^2}}$$
  
  \item \textbf{Spearman $\rho$:} Measures monotonic (rank-based) correlation --- captures non-linear but consistent trends.
\end{itemize}

Both methods return a correlation coefficient ($\in [-1, 1]$) and a p-value indicating statistical significance. With real TT data, these coefficients reflect \emph{actual} network-to-business relationships in the Tunisian market --- a finding of publishable quality.

\subsection{Output}

Each pipeline run produces \textbf{10 correlation results} (5 pairs $\times$ 2 methods), persisted to the \texttt{correlation\_insights} table with full traceability (run\_id, time window, p-values).

\section{Model Training on Real Data}

\subsection{Retraining Strategy}

The transition from synthetic to real data follows a phased approach:

\begin{enumerate}
  \item \textbf{Phase 1 (current):} Models are trained with v2.0 on synthetic data at $N = 3000$ samples. This establishes the architecture and inference pipeline.
  \item \textbf{Phase 2 (pending real data):} Models are retrained as v3.0 on real TT data. The number of training samples, time windows, and anomaly prevalence will be determined by the actual dataset size. Key adjustments:
  \begin{itemize}
    \item IsolationForest contamination reduced to $[0.01, 0.03]$
    \item GBR risk labels computed from real KPI observations
    \item BSS model expanded to 7 features (adding APPU, DOU)
  \end{itemize}
\end{enumerate}

\subsection{Model Persistence and Versioning}

Trained models are serialised as \texttt{.joblib} files and stored on a Docker volume (\texttt{/app/models/}). The AI service loads existing models on startup; retraining occurs only when models are absent or when a version bump is triggered.

Model versions:
\begin{itemize}
  \item \textbf{v2.0:} Synthetic-trained (current, operational)
  \item \textbf{v3.0:} Real-data-trained (pending data delivery)
\end{itemize}

\section{Evaluation Framework}

\subsection{SLA Risk Model (Supervised)}

\begin{itemize}
  \item Train/validation/test split (e.g., 70/15/15 or time-based)
  \item Metrics: MAE, RMSE, $R^2$ on held-out test set
  \item Feature importance analysis for explainability
\end{itemize}

\subsection{Anomaly Models (Unsupervised)}

When labelled validation data is available (known incidents from TT):
\begin{itemize}
  \item \textbf{Precision:} Of all records flagged as anomalies, what fraction are true anomalies?
  \item \textbf{Recall:} Of all true anomalies, what fraction were detected?
  \item \textbf{F1 Score:} Harmonic mean of precision and recall.
  \item \textbf{Confusion matrix:} True positives, false positives, true negatives, false negatives.
\end{itemize}

When labels are unavailable:
\begin{itemize}
  \item Inject known synthetic anomalies into real data and measure detection rate.
  \item Monitor false positive rate on known-stable periods.
  \item Use anomaly score distribution analysis (bimodal distribution implies good separation).
\end{itemize}

\subsection{Correlation Engine}

\begin{itemize}
  \item Statistical significance: p-value $< 0.05$ for meaningful correlations.
  \item Comparison of Pearson vs. Spearman to identify non-linear relationships.
  \item Cross-temporal stability: correlations should be consistent across different time windows.
\end{itemize}

\section{Chapter Summary}

This chapter documented the four components of the intelligence layer: GBR for SLA risk prediction, two IsolationForest models for OSS and BSS anomaly detection, and a Pearson/Spearman correlation engine for O+B convergence. Each model's design, features, training pipeline, and evaluation framework were detailed. With real TT data, these models transition from academic demonstrations to industrial-grade analytics tools. The next chapter covers the implementation details.
